%------------------------------------------------------------------------------%
\section{Identidade de Euler}
\label{sec:identidade_Euler}

Com os devidos cuidados, os resultados sobre sequências e séries
valem também para os números ou funções complexos. Vamos apresentar aqui a
\emph{Identidade de Euler}\index{Identidade de Euler}\index{Euler!identidade}
\[
e^{i\theta} = \cos\theta + i\sin\theta   \;.
\]
Na Seção~\ref{sec:NumerosComplexos} apresentamos um resumo sobre os
números complexos, aqui vamos utilizar que um número complexo
pode ser escrito como
\[
z = a + bi
\]
onde $a$ e $b$ são números reais e $i$ é a unidade imaginária
\[
i = \sqrt{-1}   \;.
\]
As operações de soma e produto com números complexos obedecem as mesmas
propriedades conhecidas para os números reais. Aplicando essas propriedades
podemos calcular as potências da unidade imaginária, por exemplo,
\[
i^2 = -1, \qquad i^3 = -i, \qquad i^4 = 1, \qquad i^5 = i   \;.
\]
Vamos estender a definição das funções seno, cosseno e exponencial
para variáveis complexas, isso é, queremos avaliar essas funções para qualquer
valor complexo $z\in\C$. Fazemos isso substituindo o $x$ real pelo $z$ complexo
nas Séries de Taylor dessas funções
\begin{gather*}
  \sin(z) = z - \frac{z^3}{3!} + \frac{z^5}{5!} + \frac{z^7}{7!} - \cdots \\[3mm]
  \cos(z) = 1 - \frac{z^2}{2!} + \frac{z^4}{4!} + \frac{z^8}{8!} - \cdots \\[3mm]
  e^z = 1 + z + \frac{z^2}{2!} + \frac{z^3}{3!} + \frac{z^4}{4!} +  \cdots
\end{gather*}
Note que, estamos apresentando esse resultado sem a devida justificativa.
Para sermos rigorosos precisamos verificar que todas as operações estão bem
definidas e que as séries convergem, porém, essa analise foge ao escopo desse curso.

Para verificarmos a Identidade de Euler escrevemos Série de Taylor
da exponencial de $z=i\theta$, com $\theta\in\R$
\[
e^{i\theta}
= 1 + (i\theta)
+ \frac{(i\theta)^2}{2!}
+ \frac{(i\theta)^3}{3!}
+ \frac{(i\theta)^4}{4!}
+ \frac{(i\theta)^5}{5!}
+ \frac{(i\theta)^6}{6!}
+ \cdots
\]
Utilizando a propriedade de potências $(xy)^r = x^r y^r$ reescrevemos essa expressão como
\[
e^{i\theta}
= 1 + i\theta
+ \frac{i^2\theta^2}{2!}
+ \frac{i^3\theta^3}{3!}
+ \frac{i^4\theta^4}{4!}
+ \frac{i^5\theta^5}{5!}
+ \frac{i^6\theta^6}{6!}
+ \cdots
\]
Utilizando o cálculo das potencias de $i$ obtemos
\[
e^{i\theta}
= 1
+ i\theta
-  \frac{\theta^2}{2!}
- i\frac{\theta^3}{3!}
+  \frac{\theta^4}{4!}
+ i\frac{\theta^5}{5!}
-  \frac{\theta^6}{6!}
+  \cdots
\]
Nessa expressão, os termos com potências impares contém $i$ e os termos
com potências pares não. Vamos agrupar esses dois tipos de termos
e colocar $i$ em evidência
\[
e^{i\theta}
= \left(
1
- \frac{\theta^2}{2!}
+ \frac{\theta^4}{4!}
- \frac{\theta^6}{6!}
+ \cdots
\right)
+ i \left(
\theta
- \frac{\theta^3}{3!}
+ \frac{\theta^5}{5!}
- \frac{\theta^7}{7!}
+ \cdots
\right)   \;.
\]
Os termos dentro do primeiro par de parênteses
são os termos da Série de Taylor do cosseno e os termos
dentro do segundo par de parênteses são a série do seno.
Podemos, então, escrever a Identidade de Euler
\[
e^{i\theta} = \cos\theta + i\sin\theta   \;.
\]
Uma consequência interessante dessa identidade é que podemos escrever
uma relação envolvendo ``Todas as Constantes da Matemática''.
Para isso basta escolher $\theta=\pi$
\[
e^{i\pi} \;=\; \cos\pi + i\sin\pi \;=\; -1 + i0 \;=\; -1
\]
ou seja
\[
e^{i\pi} + 1 = 0   \;.
\]

